\documentclass[12pt, a4paper]{article}
\usepackage{amsmath}
\usepackage{graphicx}
\usepackage{hyperref}

\hypersetup{colorlinks=true,
	linkcolor=black,
	citecolor=black,
	urlcolor=black}

\title{Control Systems Practical 1}
\date{21 August 2018}
\author{Q. Kruger, 216008466 \and R. de Bruyn, 216969696}

\begin{document}
	\pagenumbering{gobble}
	\begin{titlepage}
		\centering
		{\LARGE Controls Systems Practical 1 \par}
		\vspace*{1.5cm}
		{\large Q. Kruger, 216008466 \par}
		{\large R. de Bruyn, 216054484 \par}
		\vspace*{1.2cm}
		{\large 21 August 2018}
		% \makebox[\textwidth]{\includegraphics[width=\textwidth]{img/UJ.jpg}}
		\vspace*{\fill}
		\includegraphics[width=\textwidth]{img/UJ.jpg}
		\vspace*{\fill}
	\end{titlepage}
	\pagenumbering{roman}
	\tableofcontents
	\newpage
	\pagenumbering{arabic}
	
	\section{Introduction} % (fold)
	\label{sec:introduction}
		This is an introductory practical to control systems computing using MATLAB/Octave. The purpose of this practical is to become familiar with the common functions used to analyze control systems, which include generating and manipulating polynomials and transfer functions, as well as performing partial-fraction expansions.
	% section introduction (end)

	\section{Equipment and Tools} % (fold)
	\label{sec:equipment_and_tools}
		The only equipment used during this practical is the MATLAB/Octave scientific computing platforms, with the appropriate packages.
	% section equipment_and_tools (end)

	\section{Prelab} % (fold)
	\label{sec:prelab}
		\subsection*{Question 1} % (fold)
		\label{sub:question_1}
			\begin{itemize}
				\item[a)]
					\[
						\begin{array}{rcl}
							s^3 + 9s^2 + 26s + 24 & = & (s+2)(s^2 + 7s + 12) \\
							& = & (s+2)(s+3)(s+4)
						\end{array}
					\]
					Therefore $s_{roots} = \{-2,-3,-4\}$
				\item[b)]
					\[
						\begin{array}{rcl}
							s^4 + 12s^3 + 41s^2 + 30s & = & s(s^3 + 12s^2 + 41s +30) \\
							& = & s(s+1)(s^2 + 11s + 30) \\
							& = & s(s+1)(s+5)(s+6)
						\end{array}
					\]
					Therefore $s_{roots} = \{0,-1,-5,-6\}$
				\item[c)] $P_3 = s^4 + 13s^3 + 50s^2 + 56s + 24$\\
						  $P_4 = -s^4 -11s^3 -32s^2 -4s + 24$ \\
						  $P_5 = s^7 + 21s^6 + 175s^5 +735s^4 + 1624s^3 + 1764s^2 + 720s$
			\end{itemize}
		% subsection question_1 (end)

		\subsection*{Question 2} % (fold)
		\label{sub:question_2}
			\begin{equation*}
				\begin{array}{rcl}
					P_6 & = & (s+7)(s+8)(s+3)(s+5)(s+9)(s+10) \\
					& = & (s^2+15s+56)(s^2+8s+15)(s^2+19s+90) \\
					& = & (s^4 + 23s^3 +191s^2+743s+840)(s^2+19s+90) \\
					& = & s^6 + 42s^5 +718s^4 + 6\,732s^3 + 30\,820s^2 + 76\,530s + 7\,560
				\end{array}
			\end{equation*}
		% subsection question_2 (end)

		\subsection*{Question 3a}
		\label{sub:question_3a}
			\begin{equation*}
				\begin{array}{rcl}
					G_1(s) & = & \frac{20(s+2)(s+3)(s+6)(s+8)}{s(s+7)(s+9)(s+10)(s+15)}\\
					& = & \frac{20(s^3+5s^2+6s+6s^2+30s+36)(s+8)}{(s^3+7s^2+9s^2+63s)(s+10)(s+15)}\\
					& = & \frac{20(s^3+41s^2+36s+36)(s+8)}{(s^3+16s^2+63s)(s+10)(s+15)}\\
					& = & \frac{20s^4+380s^3+2480s^2+6480s+5760}{s^5+26s^4+223s^3+630s^2+15s^4+390s^3+3345s^2+9450s}\\
					& = & \frac{20s^4+380s^3+2480s^2+6480s+5760}{s^5+41s^4+613s^3+3975s^2+9450s}
				\end{array}
			\end{equation*} 

		\subsection*{Question 3b}
		\label{sub:question_3b}
			Now using the Rationale Root Theorem to find the factors of the numerator and denominator, we find that $(s+2)$ is a factor of the numerator
			\begin{equation*}
				\begin{array}{rcl}
				(s+2)(s^2+7s+12) & = & (s+2)(s+3)(s+4) \\
				\end{array}
			\end{equation*}
			The denominator can be given as
			\begin{equation*}
				\begin{array}{rcl}
				(s^4+12s^3++41s^2+30s) & = & (s)(s^3+12s^2+41s+30) \\
				\end{array}
			\end{equation*}
			Applying the same theorem as before, the second term in the equation* has a common factor given as $(s+1)$. The denominator is then given as 
			\begin{equation*}
					s(s+1)(s+6)(s+5)
			\end{equation*}
			Thus the final equation is given as 
			\begin{equation*}
				\begin{array}{rcl}
				G_2(s) & = & \frac{(s+2)(s+3)(s+7)}{s(s+1)(s+6)(s+5)}
				\end{array}
			\end{equation*}


		\subsection*{Question 3c}
		\label{sub:question_3c}
		\begin{equation*}
			\begin{array}{rcl}
				G_3(s) & = & G_1(s)+G_2(s)\\
				& = & \frac{20(s+2)(s+3)(s+6)(s+8)}{s(s+7)(s+9)(s+10)(s+15)} + \frac{(s+2)(s+3)(s+4)}{s(s+1)(s+6)(s+5)}\\
				& = & \frac{20(s+2)(s+3)(s+6)(s+8)(s+1)(s+5)+(s+2)(s+3)(s+4)(s+7)(s+9)(s+10)(s+15)}{s(s+7)(s+9)(s+10)(s+15)(s+1)(s+6)(s+5)}\\
			\end{array}
		\end{equation*}
		\begin{equation*}
			\begin{array}{rcl}
				G_4(s) & = & G_1(s)-G_2(s)\\
				& = & \frac{20(s+2)(s+3)(s+6)(s+8)}{s(s+7)(s+9)(s+10)(s+15)} - \frac{(s+2)(s+3)(s+4)}{s(s+1)(s+6)(s+5)}\\
				& = & \frac{20(s+2)(s+3)(s+6)(s+8)(s+1)(s+5)-(s+2)(s+3)(s+4)(s+7)(s+9)(s+10)(s+15)}{s(s+7)(s+9)(s+10)(s+15)(s+1)(s+6)(s+5)}\\
			\end{array}
		\end{equation*}
		\begin{equation*}
			\begin{array}{rcl}
				G_5(s) & = & G_1(s)G_2(s)\\
				& = & \frac{20(s+2)(s+3)(s+6)(s+8)}{s(s+7)(s+9)(s+10)(s+15)}{} \times \frac{(s+2)(s+3)(s+7)}{s(s+1)(s+6)(s+5)}\\
				& = & \frac{20(s+2)^2(s+3)^3(s+8)}{s^2(s+1)(s+5)(s+9)(s+10)(s+15)}
			\end{array}
		\end{equation*}

	\subsection*{Question 4a}
		\label{sub:question_4a}
		\begin{equation*}
			\begin{array}{rcl}
				G_6(s) & = & \frac{5(s+2)}{s(s^2+8s+15)}\\
				& = & \frac{5(s+2)}{s(s+5)(s+3)}\\
				& = & \frac{K_1}{s} + \frac{K_2}{s+5} + \frac{K_1}{s+3}
			\end{array}
		\end{equation*}

		\noindent Using Heaviside's method, we get

		\begin{equation*}
			\begin{array}{rcl}
				K_1 & = & \frac{5(s+2)}{(s+5)(s+3)} \bigg\rvert_{s=0} = \frac{2}{3}\\
				\\
				K_2 & = & \frac{5(s+2)}{s(s+3)} \bigg\rvert_{s=-5} = \frac{-3}{2} \\
				\\
				K_3 & = & \frac{5(s+2)}{(s+5)} \bigg\rvert_{s=-3} = \frac{5}{6} \\
			\end{array}
		\end{equation*}
		And thus we can express $G_6(s)$ as
		\[
			G_6(s) = \frac{\frac{2}{3}}{s} - \frac{\frac{2}{3}}{s+5} + \frac{\frac{5}{6}}{s+3}\\
		\]

	\subsection*{Question 4b}
		\label{sub:question_4b}
		\[
			\begin{array}{rcl}
				G_7(s) & = & \frac{5(s+2)}{s(s^2 + 6s + 9)}\\
				& = & \frac{5s+10}{s(s+3)^2}
			\end{array}
		\]
		\noindent Expanding $G_7(s)$ into partial fractions, we get the equation
		\[
			\frac{5s+10}{s(s+3)^2} = \frac{A}{s} + \frac{B}{(s+3)} + \frac{C}{(s+3)^2}
		\]

		\noindent We can then get the coefficients by using Heaviside's method:
		\[
			\begin{array}{rcl}
				A & = & \frac{5s + 10}{(s+3)^2}\bigg\rvert_{s=0}\\
				& = & \frac{10}{9}
			\end{array}
		\]
		\[
			\begin{array}{rcl}
				B & = & \frac{d}{ds}\left(\frac{5s + 10}{s}\right)\bigg\rvert_{s=-3}\\
				& = & \frac{5s - 5s - 10}{s^2}\bigg\rvert_{s=-3}\\
				& = & -\frac{10}{9}
			\end{array}
		\]
		\[
			\begin{array}{rcl}
				C & = & \frac{5s + 10}{s}\bigg\rvert_{s=-3} \\
				& = & \frac{5}{3}
			\end{array}
		\]

		\noindent Therefore we have
		\[
			G_7(s) = \frac{10}{9} \cdot \frac{1}{s} - \frac{10}{9} \cdot \frac{1}{s+3} + \frac{5}{3} \cdot \frac{1}{(s+3)^2}
		\]
	\subsection*{Question 4c}
	\label{sub:question_4c}
		\begin{equation*}
			\begin{array}{rcl}
				G_8 (s) & = & \frac{s(s+2)}{s(s^2+6s+34)} \\
				& = & \frac{A}{s}+\frac{B}{s+3-j5} + \frac{C}{s+3+j5} \\ 
			\end{array}
		\end{equation*}

		\noindent We can then find all the constants using Heaviside's Method:
		\begin{equation*}
			\begin{array}{rcl}
				A  & = & \frac{5(s+2)}{(s^2+6s+34)}\bigg\rvert_{s=0}\\
				& = & \frac{10}{34} \\
				& = & \frac{5}{17} 
			\end{array}
		\end{equation*}
		
		\[
			\begin{array}{rcl}
				B & = & \frac{5(s+2)}{s(s+3+j5)}\bigg\rvert_{s=-3+j5} \\
				& = & \frac{5(-1+j5)}{(-3+j5)(-3+j5+3j5)} \\
				& = & \frac{-5+j25}{-50-j30} \\
				& = & \frac{5\sqrt{26} \angle -1.3734}{10\sqrt{34} \angle 0.5405} \\
				& = & 0.4372 \angle -1.9138 \\
				& = & -0.1471 - j0.4117
			\end{array}
		\]
		
		\[
			\begin{array}{rcl}
				C & = & \frac{5(s+2)}{s(s+3-j5)}\bigg\rvert_{s=-3-j5} \\
				& = & \frac{5(-1-j5)}{(-3-j5)(-3-j5+3-j5)} \\
				& = & \frac{-5-j25}{-50+j30} \\
				& = & \frac{5\sqrt{26} \angle 1.3734}{10\sqrt{34} \angle -0.5405} \\
				& = & 0.4372 \angle 1.9138 \\
				& = & -0.1471 + j0.4117
			\end{array}
		\]

		\noindent Note that all of the values in this question have been rounded for the purposes of displaying the numbers. The high precision values were stored in the calculator's memory and used for calculations.

	% section prelab (end)


	\section{Lab} % (fold)
	\label{sec:lab}
		\subsection*{Question 1} % (fold)
		\label{sub:question}
			This question was done by using the \texttt{roots} command. The polynomials were initialized as \par
			\texttt{P\_1 = [0,1,9,26,24];} \par
			\texttt{P\_2 = [1,12,41,30,0];}

			\noindent The resulting polynomials for this questions were obtained by \par
			\texttt{P\_3 = P\_1 + P\_2;} \par
			\texttt{P\_4 = P\_1 - P\_2;} \par
			\texttt{P\_5 = conv(P\_1, P\_2);} \par

			\noindent Running the \texttt{roots} command on the above polynomials, the roots were found to be
			\begin{equation*}
				\begin{array}{rcl}
					s_{P3} & = & \{-6.72,-4.81,-0.733 + j0.452, -0.733 - j0.452\} \\
					s_{P4} & = & \{-5.27 + j0.452,-5.27 - j0.452, -1.19, 0.722\} \\
					s_{P5} & = & \{0,-1,-2,-3,-4,-5,-6\}
				\end{array}
			\end{equation*}
		% subsection question (end)

		\subsection*{Question 2} % (fold)
		\label{sub:question_2}
			\texttt{P\_6 = (tf('s')+7)*(tf('s')+8)*(tf('s')+3)*(tf('s')+5)*(tf('s')+9)*(tf('s')+10);} \par\noindent
			The result of the above command is
			\[
				P_6 = s^6 + 42 s^5 + 718 s^4 + 6372 s^3 + 30817 s^2 + 76530 s + 75600
			\]

		% From the ``symbolic'' package we can use the expand function that expands $G_1(s)$ into a polynomial.\\
		% \texttt{syms s;}\\
		% \texttt{poly = expand((s + 7)*(s + 8)*(s + 3)*(s + 5)*(s + 9)*(s + 10));}\\

		% \noindent Running the above commands gives the following result:\\
		% $s^6 + 42 s^5 + 718 s^4 + 6372 s^3 + 30817s^2+76530s+75600$
		% subsection question_2 (end)

		\subsection*{Question 3} % (fold)
		\label{sub:question_3}
			\texttt{s = tf('s');} \par\noindent
			\texttt{G\_1 = 20*(s+2)*(s+3)*(s+6)*(s+8)/(s*(s+7)*(s+9)*(s+10)*(s+15));}

			\noindent Running the above commands gives the following result:
			\[
				G_1(s) = \frac{20 s^4 + 380 s^3 + 2480 s^2 + 6480 s + 5760}{s^5 + 41 s^4 + 613 s^3 + 3975 s^2 + 9450 s}
			\]

		% subsection question_3 (end)

		\subsection*{Question 4} % (fold)
		\label{sub:question_4}
		\texttt{syms s;}\par\noindent
		\texttt{fraction = (s\^{}3+9*s\^{}2+26*s+24)/(s\^{}4+12*s\^{}3+41*s\^{}2+30*s)}\par\noindent
		\texttt{simplification = simplify(factor(fraction));}\par\noindent
		
		By printing the above result we get:
		\[
			\frac{(s+2)(s+3)(s+4)}{s(s+1)(s+5)(s+6)}
		\]


		% subsection question_4 (end)

		\subsection*	{Question 5}
		\label{sub:question_5}
			\texttt{syms s;}\par\noindent
			\texttt{result = factor(simplify(G1+G2));}\\\par\noindent
			Printing \texttt{result} to screen yields the following result:
				$(s+3)(s^5+305s^4+5237s^3+30987s^2+73110s+66600)/s(s+7)(s+9)(s+10)(s+15)(12s+15)$

		\subsection*{Question 6} % (fold)
		\label{sub:question_6}
			For this section, the residue command was used to calculate the coefficients of the transfer functions partial fraction expansion. We will show the procedure for calculating the first subquestion, $G_6(s)$.

			\noindent\texttt{num = [0,0,5,10];}\par\noindent
			\texttt{den = [1,8,15,0];}\par\noindent
			\texttt{[r,p,k,e] = residue(num,den);}

			\noindent Where the vector \texttt{r} represents the coefficient in the partial factor expansion. Printing \texttt{r} to screen yields the following results:
			\[
				r_a =
				\begin{bmatrix}
					0.66667 \\
					0.83333 \\
					-1.5
				\end{bmatrix}
			\]
			which corresponds to our previous results of $A=\tfrac{2}{3}$, $C=\tfrac{5}{6}$, $B=\tfrac{-3}{2}$, respectively.

			Doing the same procedure on the other subquestions, we get the following results:
			\[
				r_b =
				\begin{bmatrix}
					-1.1111 \\
					1.6667 \\
					1.1111
				\end{bmatrix}
			\]
			\[
				r_c =
				\begin{bmatrix}
					-0.14706 - j0.41176 \\
					-0.14706 + j0.41176 \\
					0.29412
				\end{bmatrix}
			\]

			\noindent Which corresponds to the coefficients found in \mbox{\hyperref[sub:question_4b]{question 4b}} and \mbox{\hyperref[sub:question_4c]{question 4c}}.
		% subsection question_6 (end)
	% section lab (end)sub:question_4c


	\section{Postlab} % (fold)
	\label{sec:postlab}
		Octave was able to do in a few minutes what would otherwise take a human half an hour to do by hand, and is a very convenient and intuitive way to do scientific computing. An interesting finding in the course of this practical, is that polynomials can be multiplied using convolution. This is due to the complexity of the Fast Fourier Transform, which has computational complexity $\mathcal{O}(N\log N)$. Doing the ``brute-force'' straight forward computation has complexity $\mathcal{O}(N^2)$\cite{polyconvolution}. Often times, performing the FFT, doing the multiplication, and then doing the inverse transform, is more computationally efficient when performing multiplication.
	% section postlab (end)

	\begin{thebibliography}{9}
		\bibitem{polyconvolution}
		The Polynomial Multiplication Problem, \url{http://algorithm.cs.nthu.edu.tw/~course/Extra_Info/Divide%20and%20Conquer_supplement.pdf}
	\end{thebibliography}

\end{document}